% !TeX program = lualatex
% Material complementar: documento do aluno. Compila com pdfLaTeX ou LuaLaTeX; [final] quando não houver pendências.
\documentclass[11pt]{article}
\usepackage{aula-geo}
\graphicspath{{../../../fontes/provas-recortes/}}
../estilo/metadados.tex
\definir\tipodocumento{Material complementar}
\definir\versao{v0}

\begin{document}
\blocotitulo

\section*{Como usar este material}
\pendente{o que a aula fez e como estas questões continuam o treino}

\begin{destaque}[title=Resumo]
\pendente{o essencial da aula em poucas linhas}
\end{destaque}

\section{Questões}
% \questaoenem{2013_CH_001}{2013}{1}{E}
\begin{questaoautoral}{B}
\pendente{enunciado autoral}
\alternativas
\item A \item B \item C \item D \item E
\fimalternativas
\end{questaoautoral}

\begin{redacao}
\pendente{repertório e eixos para a redação}
\end{redacao}

\section*{Gabarito}
\imprimirgabarito
\end{document}
